% , e.g. SQuAD~\cite{DBLP:conf/emnlp/RajpurkarZLL16} and NarrativeQA~\cite{narrativeqa} challenges.
% \todo{motivation for using KB for the QA task mentioning briefly some related work on QA from text}
% Natural language interfaces to structured data sources, such as Knowledge Graphs (KGs), databases and tables, make them more accessible to a wider user audience, especially the non-expert users, who do not have the skills in query languages such as SQL or SPARQL.
% It is also an essential component for conversational assistants.
% \todo[inline]{citations needed!}
% One of the core functions that has to be supported in the context of data access via natural language is question answering (QA), i.e. providing an answer to a natural-language question from a data source.
% It is a widely studied but still a very compelling problem.
% The major challenges of the QA task lie in (1) \textit{ambiguity} of natural languages, which leads to uncertainty in question interpretation; (2) \textit{knowledge gaps}, i.e. our semantic data sources may be incomplete and insufficient to interpret the question correctly.
% \todo{how do we use this definition?}

% now: defining complex questions
% \textit{Complex question} are usually defined in the literature as questions in contrast to simple questions that require more than one triple ~\cite{DBLP:conf/coling/SorokinG18}
% \todo{introduce two hops assumption as the upper-bound on the number of hops needed to retrieve the answer with references}

% by identifying and examining the relevant subset of the knowledge graph triples.

% % \todo{why is it challenging?}
% We hypothesize on the two main reasons that make question answering over semi-structured sources especially challenging: (1) \textbf{Representation mismatch}, the data can be stored in an arbitrary (schema-less) structure, which has to be aligned with the text of the input question; (2) \textbf{Incompleteness}, there is often additional background knowledge required to be able to make use of a structured data source, such as synonyms and paraphrases. 


% contribution
% We bridge the gap by evaluating different approaches harvesting knowledge stored in both structured (KG) and unstructured (text corpora) sources on the complex question answering task.

Here, we describe our approach for message passing within two hops in the KG but it can also be extended to an arbitrary number of hops in a similar manner.

% The state-of-the-art performance on this benchmark is 91.5\% accuracy, when accounted for multiple correct interpretations of the question~\cite{DBLP:conf/emnlp/PetrochukZ18}.
% This model consists of two components: (1) a conditional random fields (CRF) tagger with GloVe word embeddings for subject recognition given the text of the question, and (2) a bidirectional LSTM with FastText word embeddings for relation classification given the text of the question and the subject from the previous component.

% detail simple questions architectures



% The type of expected answer $A$ is usually restricted to a closed set, such an answer as a set of KG entities (SELECT) or its cardinality (COUNT queries).

% \todo{belief update would be good to add as well}
% We present an end-to-end model that links questions to entities and predicates of the KG and infers the answer based on these believes.
% Both components can be trained with back-propagation algorithm to learn the optimal weights for combining the evidence signals.


The set of labels we use for the QA task includes the set of entities and predicates that needs to be expanded on each of the hops: $<E_1, P_1>$ for the first hop, $<E_2, P_2>$ for the second hop, and the set of constraints on the answer type $T$.
Afterwards each question is represented as a sequence of labels: $Q=<l_1, ... , l_n>$, where $l_i \in$ \{``E1'', ``P1'', ``E2'', ``P2'', ``T''\}.
Greedily we collect the mentions into the corresponding sets and match them against the pre-computed entity and predicate catalogs to retrieve a ranked set of candidate URIs per span.
\todo{more details on catalogs and matching}
At the end of the interpretation phase question $Q$ is represented as a propositional formula in conjunctive normal form (CNF) under the interpretation $I$: $I(Q) = E_1 \wedge P_1 \wedge E_2 \wedge P_2 \wedge T$, where each of the and-clauses is a disjunction of the URIs modeling paraphrases.
\todo{example}
% Next, these assignments are used as the belief priors to derive the answer in the inference procedure.
% The goal of the inference step then is to trigger the message-passing algorithm and aggregate all entity and predicate activations resulting into the final probability distribution over entities to produce the answer set.
% We rely on the entity retrieval approach matching the pre-processed input question against all entity labels indexed within the entity catalog~\cite{DBLP:series/irs/Balog18}.

% It includes three consecutive steps:
% (1) extracting focused subgraph;
% (2) ranking predicates;
% (3) ranking entities.

% Since semantic relevance computation, e.g. using k-nearest neighbors for distributed vector representations (embeddings), is computationally expensive, the first step is designed to reduce the scope to a subgraph of the KG.
% This first step has to be run at least once on the whole KG and therefore, has to be computationally efficient (cheap).


% Next, we extract the subgraph S1, which includes all entities and predicates within 2 hops from each of the selected seed entities).

% We use this subgraph (S1) to get a subset of predicates $P_1 \subset P$, which we rank with respect to the input question, and use to further reduce our focused subgraph, $S_2 \subset S_1$.

% We assign scores to all entities within the subgraph $S_2$, and use them as belief priors to derive the answer in the inference procedure.

% We introduce the following notation: we denote 

% The goal of the inference phase is to 
% the CNF formula $Q_I = I(Q)$, which is the representation of the question $Q$ under the interpretation $I$.

\todo{describe the relaxation assumptions we employ}

The inference starts from the subgraph of the KG $S_1 \subset G$ induced by the seed set $<E_1, P_1>$, which is the union of all triples interpreted as relevant for the first hop:
$S_1 = \cup \forall <s, p, o>:  (s\in E_1 \lor o\in E_1) \land p\in P_1$.
% The subgraph is stored as a stack of sparse matrices for each of the predicates: $S=<S>$

The message-passing (MP) algorithm (Figure~\ref{alg:mp}) takes as input the initial activations for entities $E_1$ and predicates $P_1$ and the subgraph $S_1$ and outputs a single vector with the new activation values $a_1$. 
\todo{add dimensionality for all matrices E, P and S}
The new activation values are computed for all the entities in the subgraph $\forall e_i \in S_1$ as the weighted sum of the previous activations normalized across the components in $<E_1, P_1>$.



% pseudo code for the MP algorithm
\begin{algorithm}
\caption{MP algorithm}\label{alg:mp}
\begin{flushleft}
\textbf{Input:} E, P, S\\
\textbf{Output:} a
\end{flushleft}
\begin{algorithmic}[1]

\end{algorithmic}
\end{algorithm}

% the subgraph of the KG $S_1 \subset G$, which includes all entities and predicates from the sets $<E1, P1>$ identified at the interpretation phase.
If the second hop is required ($P2 \neq \emptyset$), we will propagate the activations from the first hop $A_1$ along with the new activations from $E2$ further following the same procedure as for the first hop starting from inducing the subgraph $S_2 \subset G$ with the seed entities and predicates $<A1 \cup E2, P2>$.


\todo{explain how to process different types of queries}

% in Figure~\ref{}
% All graph queries can be translated into equivalent matrix operations, which guarantee exact inference results~\cite{DBLP:conf/hpec/KepnerABBFGHKLM16}.
% For this we need to select, which edges in which order have to be traversed.

% However, instead of selecting from the set of all edges we introduce two simplifying assumptions allowing us to select from the set of predicates, i.e. types of edges, which helps to greatly reduce complexity of the model:
% (1) all edges in the KG are modeled as undirected;
% (2) all edges of the selected predicates can be traversed.

% These assumptions harm precision in those cases, where the direction of predicate edges bears semantic signal, e.g. `parent' relation, or the entities have an overlapping set of predicates, which has an affect on the inference results.
% However, we empirically show in our experiments that such assumptions hold true in the majority of cases for real-world KGs and the model based on these assumptions still provides high overall performance.

% \subsection{Learning}


% We leverage the current state-of-the-art approaches knowledge integration and representation approaches

% \subsection{Retrofitting word embeddings with external relations from KGs}

% We use an approach, called \textit{retrofitting}, which allows to combine representation models learned from various types of knowledge sources, such as different word embeddings and knowledge graphs~\cite{DBLP:conf/aaai/SpeerCH17}. 


% similar to the one applied to train the ConceptNet Numberbatch embeddings\footnote{\url{https://github.com/commonsense/conceptnet-numberbatch}}~\cite{DBLP:conf/aaai/SpeerCH17}


% \subsection{End-to-end question answering via Relational Graph Convolutional Networks}

% Relational Graph Convolutional Networks (R-GCNs)~\cite{DBLP:conf/esws/SchlichtkrullKB18} allow to train neural networks in an end-to-end fashion using a Knowledge Graph as an input.


% There was a number of new datasets developed recently to bootstrap the progress on the task of KG-QA. We evaluate our approach on two publicly available datasets with complex questions over DBpedia:

% \begin{enumerate}

% \item \textit{SimpleQuestions} (SQ) dataset\footnote{\url{https://github.com/davidgolub/SimpleQA/tree/master/datasets/SimpleQuestions}}~\cite{DBLP:journals/corr/BordesUCW15} contains 108,442 questions split into 70\% (75910) as training set, 10\% as
% validation set (10845), and the remaining 20\% as test set. All questions are ``simple'' questions. The answer for each of the questions is always a single triple (subject, relationship, object) extracted from the Freebase KG.

% \item \textit{WebQuestions} (WQ)\footnote{\url{https://nlp.stanford.edu/software/sempre/}}~\cite{DBLP:conf/emnlp/BerantCFL13} contains 5,810 questions constructed based on the Freebase KG and divided into the train (3,778 examples) and test splits (2,032 examples). 85\% of questions are ``simple'' questions~\cite{DBLP:conf/naacl/Yao15}.

% The answer was restricted entities, values, or list of entities.

% There are three types of answers in this dataset (1) a list of all the matched entities;

% \item \textit{GraphQuestions} (GQ)\footnote{\url{https://github.com/ysu1989/GraphQuestions}}~\cite{DBLP:conf/emnlp/SuSSSGYY16} contains 5,166 questions constructed based on the Freebase KG and is divided into the train and test splits.

% \item \textit{ComplexQuestions} (CQ)\footnote{\url{https://github.com/JunweiBao/MulCQA/tree/ComplexQuestions}}~\cite{DBLP:conf/coling/BaoDYZZ16} contains 2,100 ``complex'' questions, which require more than one relation to get the answer, and are constructed based on the Freebase KG and is divided into the train and test splits.

% \item \textit{ComplexWebQuestions} Version 1.1 (CWQ)\footnote{\url{https://www.tau-nlp.org/compwebq}}~\cite{DBLP:conf/naacl/TalmorB18} contains 34,689 questions constructed based on the Freebase KG and is divided into the train and test splits. State-of-the-art (SoA) performance: 34.2 precision~\cite{DBLP:journals/corr/abs-1807-09623}.

% \item (LQ) 

% \todo[inline]{add another dataset e.g. QALD-7 \url{https://project-hobbit.eu/challenges/qald2017}}

% The 4K SPARQL queries from the training split explicitly reference more than 3K DBpedia entities by their Unique Resource Identifiers (URIs).
% \todo[inline]{5,042 entities covered?}

% 3,180 SELECT queries in the training set

% The dataset also contains empty answers for which no matching triples were found in the DBpedia dump.


% 82\% of the questions in this dataset require more than one triple from the KG or an aggregation function (ASK or COUNT keywords) to produce the correct answer.
% There are questions with three types of answers in this dataset (1) a list of all the matched entities; (2) an integer, which is the count of all the matched entities (3) boolean, \textit{True} if at least one matching triple exists.
% We focus only on the first type of questions in our evaluation, which account for 3,180 questions in the training split.
% However, we could not retrieve any answer for 245 of these questions from the corresponding DBpedia dump, which resulted in 2,935 questions in the training dataset.
% \todo[inline]{how many different templates?}

% By examining the answer entity distribution in the training data we discovered a long tail of answers that appear in the dataset much more frequently than the others, such as (u'\url{http://dbpedia.org/resource/Chess}', 89), (u'\url{http://dbpedia.org/resource/Academy_Awards}', 56), (u'\url{http://dbpedia.org/resource/Buddhism}', 54)).
% We observed that the machine learning models discover this regularity and learn to select the most frequent answer for all the questions.
% Once being stuck in such local minimum, the model is not able to recover from it and discover any other meaningful patterns in the dataset.
% Therefore, we re-sampled the original training dataset to balance the frequency of answers more evenly by setting the upper frequency threshold to 3 and dropping all the extra question-answer pairs with the same answers.
% We released the new split to ensure reproducibility of our experiments and for the future use by the research community\footnote{\url{anonymised_url}}.
% It contains 1,151 question-answer pairs with the maximum answer entity frequency of 3.

% \todo[inline]{@Javi describe API access to HDT KG relations adjacency matrix 2 hops subgraphs}


% 249 of the \textit{SELECT} queries in the LC-QuAD train-split did not return any result and we considered an empty set to be the correct answer in those cases.

% , which hold facts about various entities, specifically their properties and relations to other entities.
% There are 217M unique node labels, in total, which correspond either to a URI or a literal, i.e. a non-unique string or number.

% answers
% We checked that 702 LC-QuAD questions can not be answered using this subset of DBpedia, i.e. the corresponding SPARQL queries against this KG version return empty sets.


% describe HDT API
\subsection{Implementation details}
\label{sec:implementation}
\textbf{HDT}.
\todo{Javi: details on the HDT format with references + short API description with link to the library}
%
\textbf{Inverted index}. Many entity references in the LC-QuAD questions can be handled using simple string similarity matching techniques, such as `companies' can be mapped to ``http://dbpedia.org/ontology/Company''.
We built an ElasticSearch (Lucene) index to efficiently retrieve such entity candidates via string similarity to their labels.
% The entity labels were automatically generated from entity URIs by stripping the domain part of the URI and lower-casing, e.g. entity ``http://dbpedia.org/ontology/Company'' received label ``company'', to be able to better match question words.

\textbf{Embeddings}. LC-QuAD questions also contain more complex paraphrases of the entity URIs that require semantic similarity computation beyond fuzzy string matching, such as ``movie'' refers to ``http://dbpedia.org/ontology/Film'', ``stockholder'' to ``http://dbpedia.org/property/owner'' or ``has kids'' to `http://dbpedia.org/ontology/child'.
We embed the entity and predicate labels with FastText~\cite{} to detect semantic similarities beyond string matching.

% fasttext
% textual similarityy

% text word semantic string

% convert magnitude annoy index increase performance 






% We used pre-trained KGlove embeddings~\cite{DBLP:conf/semweb/CochezRPP17} that encode a subset of the DBpedia graph structure (2016-04 dump) as a set of entity embeddings (200-dimensional vectors) for almost 9M DBpedia entities.
% We estimate the overlap between this representation and the entities in LC-QuAD dataset by considering the seed entities and relations, as well as the entities from the SPARQL queries and their answers.

% 8,876,676.
% \todo[inline]{check that all relations are also encoded, explain the algorithm variation}


% \textbf{Frequencies}.
% We use word and entity counts as priors for candidate selection.
% The word counts are based on the SubtlexUS project\footnote{\url{http://subtlexus.lexique.org}}~\cite{brysbaert2009moving}, which provides English word frequencies counted on a corpus of 50 million words from film and television subtitles.
% These counts reportedly outperformed Google's Books word counts for the word processing tasks by humans~\cite{brysbaert2011assessing}.



% \todo[inline]{which evaluation metric do we report}
% hits at top1 top5 or QA accuracy, which is the #correct/#total answer-question pairss
% \todo[inline]{how do we assess answers to questions that have multiple correct answers?}

% \item \textit{DBNQA} (DQ) dataset\footnote{\url{https://figshare.com/articles/Question-NSpM_SPARQL_dataset_EN_/6118505}}~\cite{hartmanngenerating} contains 894,499 questions over the DBpedia.


% \item WikiTableQuestions\footnote{\url{https://nlp.stanford.edu/software/sempre/wikitable/}}~\cite{DBLP:conf/acl/PasupatL15} contains simple as well as complex question-answer pairs that require multi-step reasoning.

% \item Sequential Question Answering (SQA) dataset\footnote{\url{https://www.microsoft.com/en-us/download/details.aspx?id=54253}}~\cite{DBLP:journals/corr/IyyerYC16} contains 17,553 questions.

% \end{enumerate}

In our experiments we used $L={'E1', 'P1', 'E2', 'P2', 'C'}$, where 'E1' and 'P1' labels mark the tokens for entity and predicate mentions that should be expanded in the first message-passing step, 'E2' and 'P2' will be expanded next together with the activations from the previous step.


% We also use the following pre-trained word embedding data sets:
% \begin{enumerate}
%   \item \label{data:embeddings_word_glove} GloVe pre-trained word embeddings ...
% \end{enumerate}

% \subsection{Set up}

% \textbf{Baseline}

% Our first baseline consists of a neural network which attempts to predict the embedded vector representation of the answer, given the question.
% In detail, the setup is as follows:
% \begin{enumerate}
%   \item We split the question into words and map each word to its word embedding, taken from the pretrained word vectors described in \cref{sec:datasets} \cref{data:embeddings_word_glove}. For each question this results in a sequence of vectors of dimension \todo[inline]{which dim}.
%   \item This sequence is then fed to a stack of three GRU layers. 
%   The whole output sequence of the first two layers is forwarded into the next layer. 
%   From the last layer only the final output is retained and passes on to the next layer.
%   \item Next, we use a dropout layer to avoid over-fitting.
%   \item The output of the dropout layer is the final output of the network.
% \end{enumerate}

% The network is trained with the question answer pairs of \todo[inline]{MC which dataset in the end}, where the question is a list of words and the answer is \todo[inline]{MC one or more}
% an entity from the knowledge graph.
% The output this network tries to predict are the embedded vectors computed from the KG, as described in 
% \todo[inline]{MC fill in which exact embedding was used}
% \cref{sec:dataset}.

% \todo[inline]{MC Fill all parameters from final experiments}
% The training parameters are as follows: batch size: 32, epochs :20, and a learning rate of 0.001 for the Adam optimizer.
% For training, we use 30\% of the data as a validation split.
% As a loss function we use the cosine distance.

% \todo[inline]{MC Likely this evaluation should be moved up. The evaluation should not only be like this for this particular experiment.}
% To evaluate the model, we use our test questions and obtain an output vector for each of them.
% The output of the network is a position in the embedding space of the KG.
% Hence, to find out which entity was predicted we have to find which entity embedding corresponds to the predicted one.

% Now, we cannot expect that the network predicts the embedding of an entity exactly.
% Hence, we look for the predicted entity by searching for the nearest neighbor(s) of the predicted vector.

% \subsection{Experiments}

% \textbf{Assumptions}. We test the base model built from the original SPARQL queries to determine the fraction of the dataset affected by our model assumptions , which shows the upper-bound performance for our approach.
% As a reminder, the simplifying assumptions are (1) undirectionality of predicates and (2) disjoint predicate sets.
% For a random sample of 453 training examples from LC-QUAD we get the following results: P 0.98 R 1.00 F 0.98

% KG entity candidate relevance ranking 
% \textbf{Entity Linking.} 

% ES index Lucene BM25



% \todo[inline]{n} mentions in the LC-QuAD questions can be efficiently resolved using only string similarity provided by the BM25 weighting scheme~\cite{}.
% We used the implementation of BM25 in Galago (The Lemur Project)\footnote{\url{https://www.lemurproject.org/galago.php}}~\cite{}.


% There was no official state-of-the-art performance for question answering on this dataset reported to date\footnote{\url{http://lc-quad.sda.tech}}.
% Accuracy reported on this dataset for the entity and relation linking tasks are 0.65 and 0.36, respectively (EARL~\cite{DBLP:conf/semweb/DubeyBCL18}).
% Other results on this benchmark were reported for the query graph ranking task alone assuming that all entities were already correctly linked.

% \todo[inline]{Our approach ...}

% A recently proposed approach~\cite{maheshwari2018learning} casts the question-answering as a subgraph ranking task.

% \textbf{Entity Linking.} To make entity linking scalable, it is usually preceded by a candidate selection step~\cite{D18-1211}, which is designed to filter a subset of entities of a bounded size, much smaller than the full entity set, to be considered as relevant matches in all the future steps.
% The standard candidate selection relies on indexing and relevance ranking of entities via BM25 weighting scheme~\cite{}.
% This approach can be efficiently implemented but it relies solely on the string-based matching and will miss all semantic matches, such as `kid' and `child'.
% One way to deal with it is semantic enrichment expanding the entity labels with additional semantic relations from external sources such as synonyms and paraphrases~\cite{DBLP:conf/semweb/DubeyBCL18}.

% \todo[inline]{describe the drawbacks of this approach}
% related work on LC-QuAD or similar datasets for complex KBQA 


% Many state-of-the-art approaches to answering simple questions over knowledge graphs are often based on the neural network architectures~\cite{DBLP:conf/www/LukovnikovFLA17,DBLP:journals/corr/abs-1804-08798}, which are tailored to answer only simple questions and can not be . 

% The current state-of-the-art approach~\cite{DBLP:conf/acl/YuYHSXZ17} that works well for both complex and simple questions reports the performance of 78.7 accuracy on the SimpleQuestions and 63.9 on the WebQuestions baselines but requires a complex pipeline-based architecture, which integrates separate entity linking, relation detection and constraint detection components.

% \subsection{Question answering from tables}

% WikiTableQuestions\footnote{\url{https://nlp.stanford.edu/software/sempre/wikitable/}}~\cite{DBLP:conf/acl/PasupatL15} provides a challenging benchmark with the current state-of-the-art of 43-46\% accuracy on the test split~\cite{DBLP:conf/emnlp/KrishnamurthyDG17}. This dataset includes simple as well as complex question-answer pairs that require multi-step reasoning.



% The model is trained end-to-end with weak supervision of question-answer pairs, and does not require domain-specific grammars, rules, or annotations that are key elements in previous approaches to program induction. The main experimental result in this paper is that a single Neural Programmer model achieves 34.2% accuracy using only 10,000 examples with weak supervision. An ensemble of 15 models, with a trivial combination technique, achieves 37.7% accuracy, which is competitive to the current state-of-the-art accuracy of 37.1% obtained by a traditional natural language semantic parser.


% Existing approaches that can be applied to complex  
% \cite{DBLP:journals/corr/abs-1803-00832}



% \subsection{Knowledge Integration Approaches}

% Word embeddings, such as word2vec~\cite{DBLP:journals/corr/abs-1301-3781}, provide a compact representation of word co-occurrences from a large text corpus, which reflect semantic relations between the words. Knowledge graph (KG) embeddings, such as RDF2Vec~\cite{DBLP:conf/semweb/CochezRPP17}, provide a compact representation of the KG structure and semantic relations between the KG entities. 

% Retrofitting
% A number of approaches were proposed for integration of knowledge from both text and KGs, such as

% \begin{enumerate}
% \item Initializing KG-entity embeddings with their word embeddings~\cite{DBLP:conf/nips/SocherCMN13}
% \item Concatenating word and KG-entity embeddings~\cite{DBLP:conf/semweb/ThomaRB17}
% \item 
% \end{enumerate}

% Finally, ensemble learning refers to an approach, when several models are trained separately and their results are integrated to produce the final answer. This approach allows to integrate very different models trained on different data sources or based on completely different approaches. Ensembles often provide superior results in practice in comparison to single models, e.g. for question answering on text (SQuAD challenge~\cite{DBLP:conf/emnlp/RajpurkarZLL16}).

% Neither of the approaches proposed above was evaluated for integrating structured (KGs) and unstructured knowledge (text) for the complex question answering task.

% \cite{} showed that an ensemble of word- and character-level language models outperform 



% Much attention of the research community today is directed towards reading comprehension and question answering (QA) from text, e.g. SQuAD~\cite{DBLP:conf/emnlp/RajpurkarZLL16} and NarrativeQA~\cite{narrativeqa} challenges, while enabling efficient information access to a variety of semi-structured data sources, such as tables and Knowledge Graphs (KGs), remains a largely under-explored area.

% % \todo{why is it challenging?}
% We hypothesize on the two main reasons that make question answering over semi-structured sources especially challenging: (1) \textbf{Representation mismatch}, the data can be stored in an arbitrary (schema-less) structure, which has to be aligned with the text of the input question; (2) \textbf{Incompleteness}, there is often additional background knowledge required to be able to make use of a structured data source, such as synonyms and paraphrases. 


% contribution
% We bridge the gap by evaluating different approaches harvesting knowledge stored in both structured (KG) and unstructured (text corpora) sources on the complex question answering task.


% \todo[inline]{Our approach ...}

% A recently proposed approach~\cite{maheshwari2018learning} casts the question-answering as a subgraph ranking task.

% \textbf{Entity Linking.} To make entity linking scalable, it is usually preceded by a candidate selection step~\cite{D18-1211}, which is designed to filter a subset of entities of a bounded size, much smaller than the full entity set, to be considered as relevant matches in all the future steps.
% The standard candidate selection relies on indexing and relevance ranking of entities via BM25 weighting scheme~\cite{}.
% This approach can be efficiently implemented but it relies solely on the string-based matching and will miss all semantic matches, such as `kid' and `child'.
% One way to deal with it is semantic enrichment expanding the entity labels with additional semantic relations from external sources such as synonyms and paraphrases~\cite{DBLP:conf/semweb/DubeyBCL18}.

% \todo[inline]{describe the drawbacks of this approach}
% related work on LC-QuAD or similar datasets for complex KBQA 


% Many state-of-the-art approaches to answering simple questions over knowledge graphs are often based on the neural network architectures~\cite{DBLP:conf/www/LukovnikovFLA17,DBLP:journals/corr/abs-1804-08798}, which are tailored to answer only simple questions and can not be . 

% The current state-of-the-art approach~\cite{DBLP:conf/acl/YuYHSXZ17} that works well for both complex and simple questions reports the performance of 78.7 accuracy on the SimpleQuestions and 63.9 on the WebQuestions baselines but requires a complex pipeline-based architecture, which integrates separate entity linking, relation detection and constraint detection components.

% \subsection{Question answering from tables}

% WikiTableQuestions\footnote{\url{https://nlp.stanford.edu/software/sempre/wikitable/}}~\cite{DBLP:conf/acl/PasupatL15} provides a challenging benchmark with the current state-of-the-art of 43-46\% accuracy on the test split~\cite{DBLP:conf/emnlp/KrishnamurthyDG17}. This dataset includes simple as well as complex question-answer pairs that require multi-step reasoning.



% The model is trained end-to-end with weak supervision of question-answer pairs, and does not require domain-specific grammars, rules, or annotations that are key elements in previous approaches to program induction. The main experimental result in this paper is that a single Neural Programmer model achieves 34.2% accuracy using only 10,000 examples with weak supervision. An ensemble of 15 models, with a trivial combination technique, achieves 37.7% accuracy, which is competitive to the current state-of-the-art accuracy of 37.1% obtained by a traditional natural language semantic parser.


% Existing approaches that can be applied to complex  
% \cite{DBLP:journals/corr/abs-1803-00832}



% \subsection{Knowledge Integration Approaches}

% Word embeddings, such as word2vec~\cite{DBLP:journals/corr/abs-1301-3781}, provide a compact representation of word co-occurrences from a large text corpus, which reflect semantic relations between the words. Knowledge graph (KG) embeddings, such as RDF2Vec~\cite{DBLP:conf/semweb/CochezRPP17}, provide a compact representation of the KG structure and semantic relations between the KG entities. 

% Retrofitting
% A number of approaches were proposed for integration of knowledge from both text and KGs, such as

% \begin{enumerate}
% \item Initializing KG-entity embeddings with their word embeddings~\cite{DBLP:conf/nips/SocherCMN13}
% \item Concatenating word and KG-entity embeddings~\cite{DBLP:conf/semweb/ThomaRB17}
% \item 
% \end{enumerate}

% Finally, ensemble learning refers to an approach, when several models are trained separately and their results are integrated to produce the final answer. This approach allows to integrate very different models trained on different data sources or based on completely different approaches. Ensembles often provide superior results in practice in comparison to single models, e.g. for question answering on text (SQuAD challenge~\cite{DBLP:conf/emnlp/RajpurkarZLL16}).

% Neither of the approaches proposed above was evaluated for integrating structured (KGs) and unstructured knowledge (text) for the complex question answering task.

% \cite{} showed that an ensemble of word- and character-level language models outperform 




% We leverage the current state-of-the-art approaches knowledge integration and representation approaches

% \subsection{Retrofitting word embeddings with external relations from KGs}

% We use an approach, called \textit{retrofitting}, which allows to combine representation models learned from various types of knowledge sources, such as different word embeddings and knowledge graphs~\cite{DBLP:conf/aaai/SpeerCH17}. 


% similar to the one applied to train the ConceptNet Numberbatch embeddings\footnote{\url{https://github.com/commonsense/conceptnet-numberbatch}}~\cite{DBLP:conf/aaai/SpeerCH17}


% \subsection{End-to-end question answering via Relational Graph Convolutional Networks}

% Relational Graph Convolutional Networks (R-GCNs)~\cite{DBLP:conf/esws/SchlichtkrullKB18} allow to train neural networks in an end-to-end fashion using a Knowledge Graph as an input.

% \section{Evaluation}

% \subsection{Datasets}
% \label{sec:datasets}
% There was a number of new datasets developed recently to bootstrap the progress on the task of KG-QA. We evaluate our approach on two publicly available datasets with complex questions over DBpedia:

% \begin{enumerate}

% \item \textit{SimpleQuestions} (SQ) dataset\footnote{\url{https://github.com/davidgolub/SimpleQA/tree/master/datasets/SimpleQuestions}}~\cite{DBLP:journals/corr/BordesUCW15} contains 108,442 questions split into 70\% (75910) as training set, 10\% as
% validation set (10845), and the remaining 20\% as test set. All questions are ``simple'' questions. The answer for each of the questions is always a single triple (subject, relationship, object) extracted from the Freebase KG.

% \item \textit{WebQuestions} (WQ)\footnote{\url{https://nlp.stanford.edu/software/sempre/}}~\cite{DBLP:conf/emnlp/BerantCFL13} contains 5,810 questions constructed based on the Freebase KG and divided into the train (3,778 examples) and test splits (2,032 examples). 85\% of questions are ``simple'' questions~\cite{DBLP:conf/naacl/Yao15}.

% The answer was restricted entities, values, or list of entities.

% There are three types of answers in this dataset (1) a list of all the matched entities;

% \item \textit{GraphQuestions} (GQ)\footnote{\url{https://github.com/ysu1989/GraphQuestions}}~\cite{DBLP:conf/emnlp/SuSSSGYY16} contains 5,166 questions constructed based on the Freebase KG and is divided into the train and test splits.

% \item \textit{ComplexQuestions} (CQ)\footnote{\url{https://github.com/JunweiBao/MulCQA/tree/ComplexQuestions}}~\cite{DBLP:conf/coling/BaoDYZZ16} contains 2,100 ``complex'' questions, which require more than one relation to get the answer, and are constructed based on the Freebase KG and is divided into the train and test splits.

% \item \textit{ComplexWebQuestions} Version 1.1 (CWQ)\footnote{\url{https://www.tau-nlp.org/compwebq}}~\cite{DBLP:conf/naacl/TalmorB18} contains 34,689 questions constructed based on the Freebase KG and is divided into the train and test splits. State-of-the-art (SoA) performance: 34.2 precision~\cite{DBLP:journals/corr/abs-1807-09623}.

% \item (LQ) 

% \textbf{Questions}.
% % \todo[inline]{add another dataset e.g. QALD-7 \url{https://project-hobbit.eu/challenges/qald2017}}
% We used LC-QuAD dataset\footnote{\url{https://github.com/AskNowQA/LC-QuAD}}~\cite{DBLP:conf/semweb/TrivediMDL17} for the evaluation of our approach.
% This dataset contains 5K questions constructed using the DBpedia KG (2016-04 version\footnote{\url{https://wiki.dbpedia.org/dbpedia-version-2016-04}}) and is divided into the train and test splits (4K + 1K questions).
% These questions were generated using a set of SPARQL templates\footnote{\url{https://github.com/AskNowQA/LC-QuAD}} by seeding them with DBpedia entities and relations, and then paraphrased by human annotators.
% % \todo[inline]{614 or 615 predicates?}
% All these templates are fit to subgraphs with diameter of at most 2-hops, contain 1-3 entities  and 1-3 relations between them.
% % The 4K SPARQL queries from the training split explicitly reference more than 3K DBpedia entities by their Unique Resource Identifiers (URIs).
% % \todo[inline]{5,042 entities covered?}

% % 3,180 SELECT queries in the training set
% There are three types of questions with the corresponding SPARQL templates: 1) \textit{SELECT} returns the subset of entities matching the query in the KG (the answer is a set of entity URIs); 2) \textit{COUNT} returns the number of entities matching the query in the KG (the answer is of type integer); 3) \textit{ASK} checks if there is at least one matching entity in the KG (the answer is of type boolean).

% All LC-QuAD questions were also annotated with entity mention spans\footnote{https://raw.githubusercontent.com/AskNowQA/EARL} to be able to separately evaluate performance on the entity linking task~\cite{DBLP:conf/semweb/DubeyBCL18}.

% % The dataset also contains empty answers for which no matching triples were found in the DBpedia dump.


% % 82\% of the questions in this dataset require more than one triple from the KG or an aggregation function (ASK or COUNT keywords) to produce the correct answer.
% % There are questions with three types of answers in this dataset (1) a list of all the matched entities; (2) an integer, which is the count of all the matched entities (3) boolean, \textit{True} if at least one matching triple exists.
% % We focus only on the first type of questions in our evaluation, which account for 3,180 questions in the training split.
% % However, we could not retrieve any answer for 245 of these questions from the corresponding DBpedia dump, which resulted in 2,935 questions in the training dataset.
% % \todo[inline]{how many different templates?}

% % By examining the answer entity distribution in the training data we discovered a long tail of answers that appear in the dataset much more frequently than the others, such as (u'\url{http://dbpedia.org/resource/Chess}', 89), (u'\url{http://dbpedia.org/resource/Academy_Awards}', 56), (u'\url{http://dbpedia.org/resource/Buddhism}', 54)).
% % We observed that the machine learning models discover this regularity and learn to select the most frequent answer for all the questions.
% % Once being stuck in such local minimum, the model is not able to recover from it and discover any other meaningful patterns in the dataset.
% % Therefore, we re-sampled the original training dataset to balance the frequency of answers more evenly by setting the upper frequency threshold to 3 and dropping all the extra question-answer pairs with the same answers.
% % We released the new split to ensure reproducibility of our experiments and for the future use by the research community\footnote{\url{anonymised_url}}.
% % It contains 1,151 question-answer pairs with the maximum answer entity frequency of 3.

% \textbf{Knowledge graph}. In our experiments we used only the English subset of the official DBpedia 2016-04 dump\footnote{\url{https://wiki.dbpedia.org/dbpedia-version-2016-04}} as the KG for retrieving answers to the LC-QuAD questions.
% Specifically, we used the publicly available version of this KG losslessly compressed into a single HDT file\footnote{http://fragments.dbpedia.org/hdt/dbpedia2016-04en.hdt}~\cite{FMPGPA:13}.

% % KG stats
% This DBpedia version contains 1B triples, 26M entities and 68,687 predicates.
% We retrieved answers for LC-QuAD questions using this KG by running the corresponding queries against our SPARQL endpoint and further used them for training our models.
% We retained only the DBpedia URIs as the correct answers to the SELECT queries, excluding all literals and other URLs.


% % \todo[inline]{@Javi describe API access to HDT KG relations adjacency matrix 2 hops subgraphs}


% % 249 of the \textit{SELECT} queries in the LC-QuAD train-split did not return any result and we considered an empty set to be the correct answer in those cases.

% % , which hold facts about various entities, specifically their properties and relations to other entities.
% % There are 217M unique node labels, in total, which correspond either to a URI or a literal, i.e. a non-unique string or number.

% % answers
% % We checked that 702 LC-QuAD questions can not be answered using this subset of DBpedia, i.e. the corresponding SPARQL queries against this KG version return empty sets.
% \subsection{Implementation}
% \label{sec:implementation}
% \textbf{HDT}.
% %
% \textbf{Inverted index}. Many entity references in the LC-QuAD questions can be handled using simple string similarity matching techniques, such as `companies' can be mapped to ``http://dbpedia.org/ontology/Company''.
% We built an ElasticSearch (Lucene) index to efficiently retrieve such entity candidates via string similarity to their labels.
% % The entity labels were automatically generated from entity URIs by stripping the domain part of the URI and lower-casing, e.g. entity ``http://dbpedia.org/ontology/Company'' received label ``company'', to be able to better match question words.

% \textbf{Embeddings}. LC-QuAD questions also contain more complex paraphrases of the entity URIs that require semantic similarity computation beyond fuzzy string matching, such as ``movie'' refers to ``http://dbpedia.org/ontology/Film'', ``stockholder'' to ``http://dbpedia.org/property/owner'' or ``has kids'' to `http://dbpedia.org/ontology/child'.
% We embed the entity and predicate labels with FastText~\cite{} to detect semantic similarities beyond string matching.

% % fasttext
% % textual similarityy

% % text word semantic string

% % convert magnitude annoy index increase performance 






% % We used pre-trained KGlove embeddings~\cite{DBLP:conf/semweb/CochezRPP17} that encode a subset of the DBpedia graph structure (2016-04 dump) as a set of entity embeddings (200-dimensional vectors) for almost 9M DBpedia entities.
% % We estimate the overlap between this representation and the entities in LC-QuAD dataset by considering the seed entities and relations, as well as the entities from the SPARQL queries and their answers.

% % 8,876,676.
% % \todo[inline]{check that all relations are also encoded, explain the algorithm variation}


% % \textbf{Frequencies}.
% % We use word and entity counts as priors for candidate selection.
% % The word counts are based on the SubtlexUS project\footnote{\url{http://subtlexus.lexique.org}}~\cite{brysbaert2009moving}, which provides English word frequencies counted on a corpus of 50 million words from film and television subtitles.
% % These counts reportedly outperformed Google's Books word counts for the word processing tasks by humans~\cite{brysbaert2011assessing}.


% \subsection{Metrics}
% In the following evaluation we report the standard performance measures for the QA task: precision (P), recall (R) and f-measure (F).
% % \todo[inline]{which evaluation metric do we report}
% % hits at top1 top5 or QA accuracy, which is the #correct/#total answer-question pairss
% % \todo[inline]{how do we assess answers to questions that have multiple correct answers?}

% % \item \textit{DBNQA} (DQ) dataset\footnote{\url{https://figshare.com/articles/Question-NSpM_SPARQL_dataset_EN_/6118505}}~\cite{hartmanngenerating} contains 894,499 questions over the DBpedia.


% % \item WikiTableQuestions\footnote{\url{https://nlp.stanford.edu/software/sempre/wikitable/}}~\cite{DBLP:conf/acl/PasupatL15} contains simple as well as complex question-answer pairs that require multi-step reasoning.

% % \item Sequential Question Answering (SQA) dataset\footnote{\url{https://www.microsoft.com/en-us/download/details.aspx?id=54253}}~\cite{DBLP:journals/corr/IyyerYC16} contains 17,553 questions.

% % \end{enumerate}





% % We also use the following pre-trained word embedding data sets:
% % \begin{enumerate}
% % 	\item \label{data:embeddings_word_glove} GloVe pre-trained word embeddings ...
% % \end{enumerate}

% % \subsection{Set up}

% % \textbf{Baseline}

% % Our first baseline consists of a neural network which attempts to predict the embedded vector representation of the answer, given the question.
% % In detail, the setup is as follows:
% % \begin{enumerate}
% % 	\item We split the question into words and map each word to its word embedding, taken from the pretrained word vectors described in \cref{sec:datasets} \cref{data:embeddings_word_glove}. For each question this results in a sequence of vectors of dimension \todo[inline]{which dim}.
% % 	\item This sequence is then fed to a stack of three GRU layers. 
% % 	The whole output sequence of the first two layers is forwarded into the next layer. 
% % 	From the last layer only the final output is retained and passes on to the next layer.
% % 	\item Next, we use a dropout layer to avoid over-fitting.
% % 	\item The output of the dropout layer is the final output of the network.
% % \end{enumerate}

% % The network is trained with the question answer pairs of \todo[inline]{MC which dataset in the end}, where the question is a list of words and the answer is \todo[inline]{MC one or more}
% % an entity from the knowledge graph.
% % The output this network tries to predict are the embedded vectors computed from the KG, as described in 
% % \todo[inline]{MC fill in which exact embedding was used}
% % \cref{sec:dataset}.

% % \todo[inline]{MC Fill all parameters from final experiments}
% % The training parameters are as follows: batch size: 32, epochs :20, and a learning rate of 0.001 for the Adam optimizer.
% % For training, we use 30\% of the data as a validation split.
% % As a loss function we use the cosine distance.

% % \todo[inline]{MC Likely this evaluation should be moved up. The evaluation should not only be like this for this particular experiment.}
% % To evaluate the model, we use our test questions and obtain an output vector for each of them.
% % The output of the network is a position in the embedding space of the KG.
% % Hence, to find out which entity was predicted we have to find which entity embedding corresponds to the predicted one.

% % Now, we cannot expect that the network predicts the embedding of an entity exactly.
% % Hence, we look for the predicted entity by searching for the nearest neighbor(s) of the predicted vector.

% \subsection{Experiments}

% \textbf{Assumptions}. We test the base model built from the original SPARQL queries to determine the fraction of the dataset affected by our model assumptions , which shows the upper-bound performance for our approach.
% As a reminder, the simplifying assumptions are (1) undirectionality of predicates and (2) disjoint predicate sets.
% For a random sample of 453 training examples from LC-QUAD we get the following results: P 0.98 R 1.00 F 0.98

% % KG entity candidate relevance ranking 
% % \textbf{Entity Linking.} 

% % ES index Lucene BM25



% % \todo[inline]{n} mentions in the LC-QuAD questions can be efficiently resolved using only string similarity provided by the BM25 weighting scheme~\cite{}.
% % We used the implementation of BM25 in Galago (The Lemur Project)\footnote{\url{https://www.lemurproject.org/galago.php}}~\cite{}.
